% ===== PRÉAMBULE CANONIQUE — à coller VERBATIM en tête de CHAQUE .tex =====
\documentclass[11pt,a4paper]{article}
\usepackage[utf8]{inputenc}\usepackage[T1]{fontenc}\usepackage{lmodern}
\usepackage[french]{babel}
\usepackage{amsmath,amssymb,mathtools}\usepackage{icomma}
\usepackage[version=4]{mhchem}          % \ce{...} : équations & formules
\usepackage{siunitx}\sisetup{locale=FR} % \qty{}{}, \si{}, \ang{}
\usepackage{chemfig}                    % \chemfig{...} : structures développées
\usepackage{xcolor}\usepackage{colortbl}
\usepackage{tikz}\usepackage{pgfplots}\pgfplotsset{compat=1.18}
\usepackage[most]{tcolorbox}\usepackage{enumitem}\usepackage{array,booktabs}
\usepackage[a4paper,margin=2.2cm]{geometry}
\usetikzlibrary{arrows.meta,calc,angles,quotes,positioning,patterns,backgrounds,shapes.geometric}
\definecolor{primary}{RGB}{222,49,99}\definecolor{primarydark}{RGB}{150,22,55}
\definecolor{defcol}{RGB}{0,114,178}\definecolor{methcol}{RGB}{52,92,148}
\definecolor{excol}{RGB}{0,135,90}\definecolor{attncol}{RGB}{200,40,55}
\tcbset{boxbase/.style={enhanced,breakable,boxrule=0.9pt,arc=2.5pt,
  left=3mm,right=3mm,top=2.3mm,bottom=2.3mm,fonttitle=\bfseries,coltitle=white}}
% --- boîtes TUTORIEL (noms EXACTS, ne pas renommer) ---
\newtcolorbox{cle}[1][]{boxbase,colback=defcol!6,colframe=defcol,
  title={\textsc{À retenir}\ifx\relax#1\relax\relax\else\textmd{ : \itshape #1}\fi}}
\newtcolorbox{methode}[1][]{boxbase,colback=methcol!7,colframe=methcol,sharp corners,
  title={\textsc{Méthode}\ifx\relax#1\relax\relax\else\textmd{ --- \itshape #1}\fi}}
\newtcolorbox{exemplebox}[1][]{boxbase,colback=excol!5,colframe=excol,
  title={\textsc{Exemple}\ifx\relax#1\relax\relax\else\textmd{ : \itshape #1}\fi}}
\newtcolorbox{piege}[1][]{boxbase,colback=attncol!6,colframe=attncol,
  title={\textsc{Piège}\ifx\relax#1\relax\relax\else\textmd{ : \itshape #1}\fi}}
\newtcolorbox{remarque}[1][]{boxbase,colback=black!4,colframe=black!45,
  title={\textsc{Remarque}\ifx\relax#1\relax\relax\else\textmd{ : \itshape #1}\fi}}
% --- boîtes EXERCICE / CORRIGÉ (noms EXACTS, auto-comptées) ---
\newtcolorbox[auto counter]{exo}[1][]{boxbase,colback=primary!4,colframe=primary,
  title={\textsc{Exercice}~\thetcbcounter\ifx\relax#1\relax\relax\else\textmd{ --- \itshape #1}\fi}}
\newtcolorbox[auto counter]{corrige}[1][]{boxbase,colback=excol!4,colframe=excol,
  title={\textsc{Corrigé}~\thetcbcounter\ifx\relax#1\relax\relax\else\textmd{ --- \itshape #1}\fi}}
\newcommand{\R}{\mathbb{R}}\newcommand{\N}{\mathbb{N}}\newcommand{\Z}{\mathbb{Z}}\newcommand{\ds}{\displaystyle}
% (un \usepackage{fancyhdr}+en-tête est possible ; il est ignoré côté web)
% =========================================================================
\begin{document}

\begin{center}{\large\bfseries\color{primarydark} Thème 1 — Niveau Moyen}\end{center}
\emph{Données (\si{\gram\per\mole}) :} $\ce{H}=1$, $\ce{C}=12$, $\ce{N}=14$, $\ce{O}=16$,
$\ce{Na}=23$, $\ce{S}=32$, $\ce{Ca}=40$, $\ce{Fe}=56$. \quad $N_{\mathrm{A}} = 6{,}02\times10^{23}\ \si{\per\mole}$.

\begin{exo}[chaîne masse $\to$ moles $\to$ molécules]
Détermine le nombre de molécules contenues dans :
\begin{enumerate}[label=\alph*)]
  \item $\SI{9}{\gram}$ d'eau \ce{H2O} \quad ($M=\SI{18}{\gram\per\mole}$) ;
  \item $\SI{8.8}{\gram}$ de dioxyde de carbone \ce{CO2} \quad ($M=\SI{44}{\gram\per\mole}$).
\end{enumerate}
\end{exo}

\begin{exo}[masse molaire du glucose, puis quantité de matière]
Le glucose a pour formule \ce{C6H12O6}.
\begin{enumerate}[label=\alph*)]
  \item Calcule sa masse molaire $M$.
  \item Quelle quantité de matière représente $\SI{90}{\gram}$ de glucose~?
  \item Et $\SI{18}{\gram}$ de glucose~?
\end{enumerate}
\end{exo}

\begin{exo}[masse d'un nombre d'entités donné]
Détermine la masse de :
\begin{enumerate}[label=\alph*)]
  \item $3{,}01\times10^{23}$ molécules de \ce{O2} \quad ($M=\SI{32}{\gram\per\mole}$) ;
  \item $1{,}204\times10^{24}$ atomes de fer \quad ($M=\SI{56}{\gram\per\mole}$).
\end{enumerate}
\end{exo}

\begin{exo}[pourcentage massique d'un élément]
Calcule le pourcentage massique de l'élément indiqué dans chaque composé.
\begin{enumerate}[label=\alph*)]
  \item \% d'oxygène dans \ce{H2O} ;
  \item \% d'azote dans le nitrate d'ammonium \ce{NH4NO3} \quad ($M=\SI{80}{\gram\per\mole}$) ;
  \item \% de calcium dans \ce{CaCO3} \quad ($M=\SI{100}{\gram\per\mole}$).
\end{enumerate}
\end{exo}

\begin{exo}[compter les atomes]
Réponds aux deux questions indépendantes.
\begin{enumerate}[label=\alph*)]
  \item Combien d'atomes d'oxygène (en nombre) contient $\SI{0.2}{\mole}$ de \ce{H2SO4}~?
  \item De $\SI{8}{\gram}$ de \ce{O2} ou de $\SI{8}{\gram}$ de soufre \ce{S}, quel échantillon
        contient le plus d'atomes~? ($M(\ce{S})=\SI{32}{\gram\per\mole}$.)
\end{enumerate}
\end{exo}

\end{document}
